\chapter{Null extensions and material closure}
\label{chap:realizaciones-nula-material}
\index{photon!physical interface}\index{electron!material closure}
\index{null extension}\index{causal cone partitioned}

\section{\(4\) realizations of the common material state: null, massive, composite, and topological}
\label{sec:cuatro-modos-realizacion}

An isolated digit determines a coordinate but does not identify the structure
that produces it, the memory it preserves, or the operation that converts it into
a magnitude. HMT lies prior to that projection: its primary object is a
compatible extension that preserves the arithmetic laws, TRIT orientation,
quotients, phase, and boundary memory. Dimensional Mechanics
subsequently specifies which component is realized as a magnitude, which remains
as memory, and which transformation produces the observable.

The developments of this section distinguish four modes by their operation of
closure:
\begin{enumerate}
 \item an open null extension transports a perturbation between
 boundaries;
 \item a closed and persistent extension acquires interior norm and defines
 the material candidate;
 \item a family of extensions weighted by action defines a thermal
 ensemble;
 \item a connection that compares sections at distinct points provides the
 codomain of the gravitational realization.
\end{enumerate}
The common source does not identify the four codomains. The corresponding
physical maps remain discontinuous:
\[
 \mathfrak R_\gamma:
 \Omega_{\HMT}^{\rm surv}\dashrightarrow\mathcal H_{\rm rad},
 \qquad
 \mathfrak R_m:
 \Omega_{\HMT}^{\rm surv}\dashrightarrow\mathcal H_{\rm mat},
\]
\[
 \mathfrak R_\Theta:
 \mathcal P(\Omega_{\HMT}^{\rm surv})\dashrightarrow\mathcal S_{\rm th},
 \qquad
 \mathfrak R_{\rm grav}:
 \Omega_{\HMT}^{\rm surv}\dashrightarrow\mathsf{Cartan}_{1,3}.
\]
Their unity resides in the genealogy; their difference, in the reading and in the
additional structures required by each realization.

\section{Null propagation between boundaries}
\label{sec:propagacion-nula-fotonica}

Let \(B\) be a transport basis. A \emph{null section} of sign
\(\pm\) is a map
\[
 s_\pm:B\longrightarrow\mathcal A_{\rm sp}P_\pm
\]
such that \(s_\pm(x)\in\mathcal A_{\rm sp}P_\pm\) for every \(x\in B\).
By \cref{prop:descomposicion-espectral-partida},
\[
 N_{\rm sp}(s_\pm(x))=0
 \qquad(x\in B).
\]
If \(\gamma\) is an oriented route between two boundaries,
\[
 \partial\gamma=x_{\rm abs}-x_{\rm em},
\]
the proposed photonic interface assigns to \(\gamma\) an extended null
section. The adjective \emph{extended} indicates that its identity belongs to the
complete transport and not to an isolated point.

The conjugate route is obtained by inversion,
\[
 \gamma^\vee=C_{\rm rev}\gamma,
 \qquad
 \partial\gamma^\vee=x_{\rm em}-x_{\rm abs}.
\]
Relative to the declared orientation reader, the pair
\((\gamma,\gamma^\vee)\) represents the advanced and retarded readings. The
two legs are ordered sequentially or housed in distinct fibers. They are not
automatically added within the same fiber: if the two conjugate components are
activated simultaneously, the exact identity
\[
 \boxed{N_{\rm sp}(uP_++vP_-)=uv}
\]
shows that, for \(u,v>0\), the section leaves the null boundary and enters the
positive interior.

\subsection{Length, go and return}

Once the dimensional velocity chart,
\(\widehat c=e^{-\mathcal E}\), and a positive basis \(u_L\) of the length
line have been fixed, define
\[
 L_{\rm int}(\gamma)=\widehat c\,|\gamma|\,u_L.
\]
Inversion preserves the number of steps:
\[
 L_{\rm int}(\gamma^\vee)=L_{\rm int}(\gamma),
 \qquad
 L_{\leftrightarrow}=2L_{\rm int}(\gamma).
\]
These identities are kinematic. Converting them into electromagnetic
propagation requires that the same map reproduce Gauss's law, the
gauge dynamics, null dispersion, and the reduction to two
transverse degrees.

\subsection{Entanglement entropy between the \(2\) readings}
\label{sec:entrelazamiento-ida-retorno}

The reversal of a route does not suffice to assert quantum entanglement. A
minimal formulation first requires a Hilbert completion. Let
\(\mathcal H_{\rm adv}\) and \(\mathcal H_{\rm ret}\) be two complex spaces
with orthonormal families
\(\{|i\rangle_{\rm adv}\}_{i\in I}\) and
\(\{|i\rangle_{\rm ret}\}_{i\in I}\), where \(I\) is finite. For a
probability \(p=(p_i)_{i\in I}\), define
\[
 |\Psi_p\rangle
 :=\sum_{i\in I}\sqrt{p_i}\,
 |i\rangle_{\rm adv}\otimes|i\rangle_{\rm ret}.
 \tag{E1}\label{eq:estado-entrelazado-ida-retorno}
\]

\begin{proposition}[Round-trip double-layer entropy]
\label{prop:entropia-entrelazamiento-ida-retorno}
The reduced matrices of the state
\(\rho_p=|\Psi_p\rangle\langle\Psi_p|\) are
\[
 \rho_{\rm adv}
 =\sum_{i\in I}p_i|i\rangle_{\rm adv}\langle i|,
 \qquad
 \rho_{\rm ret}
 =\sum_{i\in I}p_i|i\rangle_{\rm ret}\langle i|,
 \tag{E2}\label{eq:reducidas-ida-retorno}
\]
and possess the same von Neumann entropy:
\[
 \boxed{
 S_{\rm ent}(p)
 =-\operatorname{Tr}(\rho_{\rm adv}\log\rho_{\rm adv})
 =-\sum_{i\in I}p_i\log p_i.}
 \tag{E3}\label{eq:entropia-ida-retorno}
\]
Here the convention \(0\log0:=0\) is adopted. The state is separable
exactly when a single probability is equal to one.
\end{proposition}

\begin{proof}
Upon taking the partial trace, orthogonality cancels the terms with
\(i\ne j\). The nonzero eigenvalues of each reduced matrix are therefore
the \(p_i\), which gives \eqref{eq:entropia-ida-retorno}. The Schmidt
rank is one exactly in the deterministic case.
\end{proof}

Once the completion has been chosen, the proposition formalizes an entropy between
the advanced and retarded readings. It does not derive the probabilities \(p_i\), does not
identify those readings with the advanced and retarded propagators of a
field theory, and does not convert a classical route correlation into
entanglement. Construction of an isomorphism sending HMT memory
to bipartite physical states, preserving evolution, and fixing \(p\) belongs to
the \textsc{physical interface}. The null entropy of a deterministic
distribution constitutes the immediate negative control.

\begin{statusbox}
The proposition is
\textsc{Proved from explicit starting structure}: it assumes the
Hilbert completion, the orthonormal bases, and the distribution \(p\). The construction
of the intertwiner that physically realizes both layers is a
\textsc{physical interface}; it is not part of the proposition.
\end{statusbox}

\begin{definition}[HMT--MD Photonic Interface]
\label{def:interfaz-fotonica-hmt-md}
A realization photonic of HMT--MD is a section open sustained over
a direction null of the algebra split, transported between boundaries
conjugate and protected by a symmetry that preserves exactly the norm
null.
\end{definition}

The idempotents \(P_\pm\) are not, by themselves, physical photons nor the two
Maxwell polarizations. The complete realization requires an
intertwiner
\[
 \mathfrak R_\gamma:
 \partial_{\rm null}\mathcal C_{\rm sp}^{+}
 \dashrightarrow\mathcal H_{\rm Maxwell}^{\rm phys}
\]
that preserves the inner product, evolution, and gauge action, and that
selects exactly two transverse polarizations in \(3+1\)
dimensions.

\begin{criterion}[Photonic sector]
\label{crit:sector-fotonico}
The identification physical remains refuted if the evolution leaves
\(\partial_{\rm null}\mathcal C_{\rm sp}^{+}\), generates mass without the mechanism
of breaking declared, does not reproduces the dispersion null or not dejan exactly
two polarizations physical after the reduction of gauge.
\end{criterion}

\section{Material closure and electronic center}
\label{sec:clausura-material-electron}

The simultaneous coupling of two conjugate null directions yields, for
\(u,v>0\),
\[
 N_{\rm sp}(uP_++vP_-)=uv>0.
\]
No isolated component has nonzero norm; the inner norm
arises from their coupling. This algebraic statement does not add pre-existing
masses and does not yet identify the null directions with photons
\emph{on shell}.

   The structural classification of HMT contains a single fixed point of cellular inversion,
\[
 (5,5).
\]
The electron's relative action origin is \(v_e=0\). In the normalized
split realization, set
\[
 Z_e=P_++P_-=I,
 \qquad
 N_{\rm sp}(Z_e)=1.
\]
The point \((5,5)\), the origin \(v_e=0\), the raw signature
\((24,0,12,0,-12)\), and the transverse word \(555555\) are four distinct
objects.

\begin{definition}[Electronic closure candidate]
\label{def:candidato-clausura-electronica}
The resonant closure of conjugate electromagnetic null modes is
formalized as the program that seeks a normalized central material section
obtained by coupling two conjugate null modes, stabilized by
a route closure and its holonomy.
\end{definition}

The existence of the extension, the return, and the center of the atlas does not prove
minimal compactness or the stability of the closure. The outstanding dynamical
theorem must construct
\[
 \gamma_e=\xi_1\circ\cdots\circ\xi_N,
 \qquad
 \partial\gamma_e=0,
 \qquad
 \varepsilon(\gamma_e)=-1,
\]
where each \(\xi_j\) is locally zero, the total boundary vanishes, and
\(\varepsilon\) is the sign holonomy. Mass would be associated with the global
closure cocycle, not with a sum of locally zero masses. This program does not
reopen the central electron \((5,5)\), the atlas, or the already constructed
extension--route--record--character chain.

\begin{proposition}[Relativistic Causal Interior Criterion]
\label{prop:cierre-nulos-interior-causal}
Let \(p_1,\ldots,p_N\) be future null momenta in a Minkowski space with
positive temporal signature, \(p_i^2=0\), and let
\[
 P=\sum_{i=1}^{N}p_i.
\]
If all \(p_i\) are collinear, then \(P^2=0\). If at least two are
future and noncollinear, then \(P^2>0\), and one may define
\[
 m^2c^2=P^2.
\]
\end{proposition}

\begin{proof}
For future-directed null vectors,
\(p_i\cdot p_j\geq0\) holds, with equality precisely when they are collinear.
Since
\[
 P^2=\sum_i p_i^2+2\sum_{i<j}p_i\cdot p_j,
\]
the first sum is zero. All cross products vanish in the
collinear case; if a noncollinear pair exists, at least one is strictly
positive.
\end{proof}

Therefore, a closed route is not massive by mere topology, nor is an
open route null merely because it has a boundary. Realization requires that
the geometric closure and the sum of momenta enter jointly into the
causal interior.

\subsection{Möbius, conjugation, and posterior occupation}

Under the lifting and holonomy hypotheses of \cref{thm:dos-cuatro-pi-estadistica}, one turn reverses the orientation and two turns restore it. The Möbius band organizes particle and antiparticle as conjugate orientations of the same type of closure. Their physical identification requires a representation that reverses charge, preserves the norm, and intertwines the evolution. Möbius geometry does not by itself imply the algebra of fermionic occupation; its operator relations are introduced as operator hypotheses in the following chapter.

\begin{criterion}[Electron closure]
\label{crit:cierre-electronico}
The realization is refuted if the center \((5,5)\) does not lift to a
closed extension, if the holonomy does not select the odd sector, if the
Hamiltonian lacks a stable bound solution that is unique up to symmetry, or
if the charge, dispersion, and form factor do not agree with the electron
realization.
\end{criterion}
